\section{Proofs of the limits}
\label{sec:lower-proof}

This appendix contains the proofs deferred from Section~\ref{sec:limits}.
For \(d\ge2\), let \(\mu_d\) be the positive root of
\(d\mu^2-\mu-(d+2)=0\), as in Lemma~\ref{lem:bulk-potential}, and put
\(g_d=\log\frac{d+2}d\) and \(\omega_d=3\log\mu_d/g_d\). Then
\(\mu_d>1\) (Lemma~\ref{lem:bulk-potential}), and \(\lambda_d=1/\mu_d\)
is the positive root of \(f_d(y)=(d+2)y^2+y-d\). The function \(f_d\)
increases on \([0,\infty)\), as \(f_d'(y)=2(d+2)y+1>0\) there. We write
\(\kappa_0=21/20\) and
\[
 \hat\Psi(r,b)=\Psi(r,b)+\tfrac56
 =\bigl(r-\tfrac13\bigr)^2+4\bigl(r-\tfrac13\bigr)\bigl(b-\tfrac13\bigr)+\bigl(b-\tfrac13\bigr)^2 .
\]
Like \(\Psi\), the function \(\hat\Psi\) is symmetric and increasing in
each argument for real \(r,b\ge1\) (Lemma~\ref{lem:hankel-potential}(a)).

\subsection{Supplies}
\label{subsec:proof-recycling}

\begin{proof}[Proof of Lemma~\ref{lem:recycling}]
\emph{Stage relations.} Fix \(i\), and let \((r,b)\to(c_i,r)\) be the
legal stage of Definition~\ref{def:class}(b). For each block \(j\),
Corollary~\ref{cor:exact-stage} with
\((N,k,\tau)=(N_{ij},k_{ij},\tau_{ij})\) gives
\[
 E_{ij}(2M_{ij}-1)^{6k_{ij}}Q_{r,b}^{\otimes N_{ij}k_{ij}}\MM_{d_i^{2N_{ij}k_{ij}}}
 \restr M_{ij}^{6k_{ij}}Q_{c_i,r}^{\otimes N_{ij}k_{ij}}\MM_{M_{ij}^{2k_{ij}}} .
\]
Tensor these restrictions over \(j\), and use \((mT)(m'T')\cong mm'(TT')\),
\(\sum_jN_{ij}k_{ij}=W\) and~\eqref{eq:prelim-mm}. With
\(X_i=\prod_jE_{ij}(2M_{ij}-1)^{6k_{ij}}\) and
\(Y_i=\prod_jM_{ij}^{6k_{ij}}\), this gives
\(X_iQ_{r,b}^{\otimes W}\MM_{U_i}\restr Y_iQ_{c_i,r}^{\otimes W}\MM_{V_i}\).
Here \(Q_{r_i,b_i}\cong Q_{r,b}\) by~\eqref{eq:prelim-Q}. If
\(c'\le c\) and \(r'\le r\), then \(C(c,r)\restr C(c',r')\): set to zero
\(x_\alpha\) for \(\alpha\ge c'\), \(y_\eta\) for \(\eta\ge r'\) and
\(z_\zeta\) for \(\zeta\ge c'+r'-1\); the surviving terms are those with
\(\alpha<c'\) and \(\eta<r'\). Doing this in all six oriented factors
gives \(Q_{c,r}\restr Q_{c',r'}\). Hence a crop and possibly an exchange give
\(Q_{c_i,r}\restr Q_{r_{i+1},b_{i+1}}\), and
\begin{equation}\label{eq:class-stage}
 X_i\,Q_{r_i,b_i}^{\otimes W}\MM_{U_i}\restr
 Y_i\,Q_{r_{i+1},b_{i+1}}^{\otimes W}\MM_{V_i}.
\end{equation}

\emph{Induction.} We show, for \(i=0,\ldots,J\), that
\begin{equation}\label{eq:class-induction}
 r_*\MM_A\restr\Bigl(\prod_{l\le i}Y_l\Bigr)\Bigl(\prod_{l>i}X_l\Bigr)
 Q_{r_{i+1},b_{i+1}}^{\otimes W}\MM_{\mathfrak b_i} .
\end{equation}
Since \(r_*=s^{6W}\prod_lX_l\), and \(s^{6W}\MM_1\restr Q_{r_1,b_1}^{\otimes W}\)
over \(\Q\) by~\eqref{eq:prelim-Q}, tensoring with \(\MM_A\) gives the
case \(i=0\). Suppose~\eqref{eq:class-induction} holds for \(i-1\). Its
right side is a direct sum of \((\prod_{l<i}Y_l)(\prod_{l>i}X_l)\)
groups, each of \(X_i\) copies of
\(Q_{r_i,b_i}^{\otimes W}\MM_{\mathfrak b_{i-1}}\). Setting to zero the
variables of \(\MM_{\mathfrak b_{i-1}}\) with an index at least \(U_if_i\)
restricts it to \(\MM_{U_if_i}\cong\MM_{U_i}\MM_{f_i}\), so each group
restricts to \((X_iQ_{r_i,b_i}^{\otimes W}\MM_{U_i})\MM_{f_i}\). Apply
\eqref{eq:class-stage}, tensored with \(\MM_{f_i}\), and use
\(\MM_{V_i}\MM_{f_i}\cong\MM_{\mathfrak b_i}\). Each group becomes
\(Y_iQ_{r_{i+1},b_{i+1}}^{\otimes W}\MM_{\mathfrak b_i}\), which
proves~\eqref{eq:class-induction} for \(i\). Finally, a crop and
possibly an exchange give \(Q_{r_{J+1},b_{J+1}}\restr Q_{1,t}\cong\MM_{t^2}\)
by~\eqref{eq:prelim-Q}, so the case \(i=J\) gives
\(r_*\MM_A\restr m_*\MM_{t^{2W}\mathfrak b_J}=m_*\MM_B\). All maps are
restrictions with rational coefficients. No degeneration is composed
with another, since each stage relation is already exact.

(a) Since \(\mathfrak b_i=f_iV_i\le\mathfrak b_{i-1}V_i/U_i\), we get
\(\mathfrak b_J\le A\prod_iV_i/U_i\), that is, \(B/A\le\bar B/\bar A\). For the
source supply, induction on \(i\) gives
\(\mathfrak b_{i-1}=\prod_{l\ge i}U_l\prod_{l<i}V_l=U_if_i\), so every
inequality is an equality.

(b) The first stage is legal, so \(b\ge c_1+1\ge2\) in
Definition~\ref{def:class}(b), and \(s\ge2\). Hence
\(r_*/m_*=s^{6W}\prod_{i,j}E_{ij}\bigl((2M_{ij}-1)/M_{ij}\bigr)^{6k_{ij}}\ge s^{6W}>1\).
If \(B>A\), then \(\bar B/\bar A\ge B/A>1\) by (a), so both restrictions
are chain schemes. Their pooled exponents have the same positive
numerator \(\log(r_*/m_*)\), and \(\log(B/A)\le\log(\bar B/\bar A)\).

(c) Let \(\tilde U\) be the maximum in the hypothesis. We show by
induction on \(i=1,\ldots,J+1\) that
\(\mathfrak b_{i-1}\ge\frac A2\prod_{l<i}(V_l/U_l)\), and that
\(\mathfrak b_{i-1}\ge2JU_i\) if \(i\le J\). Suppose this holds for all
smaller \(i\). Then \(f_l\ge1\) for \(l<i\), and, with
\(y=\mathfrak b_{l-1}/U_l\), the bound \(\lfloor y\rfloor\ge y-1\) gives
\(\mathfrak b_l\ge(\mathfrak b_{l-1}V_l/U_l)(1-U_l/\mathfrak b_{l-1})
\ge(1-\frac1{2J})\mathfrak b_{l-1}V_l/U_l\) for \(l<i\).
By Bernoulli's inequality,
\((1-\frac1{2J})^{i-1}\ge1-\frac{i-1}{2J}\ge\frac12\), so
\(\mathfrak b_{i-1}\ge\frac A2\prod_{l<i}(V_l/U_l)\). If \(i\le J\), this is
at least \(2J\tilde U\prod_{l<i}(V_l/U_l)\ge2JU_i\), by the definition of
\(\tilde U\). At \(i=J+1\) it gives \(B/A\ge\bar B/(2\bar A)\).
\end{proof}

\paragraph{The floors in Remark~\ref{rem:composition}.}
Let \(n_\infty=AB/(B-A)\), as in Section~\ref{sec:seeds}, and let
\(\mathfrak b_0=n\ge2n_\infty\) and
\(\mathfrak b_{l+1}=\lfloor\mathfrak b_l/A\rfloor B\). Since
\(\lfloor y\rfloor\ge y-1\), we have
\(\mathfrak b_{l+1}\ge(B/A)\mathfrak b_l-B\), that is,
\(\mathfrak b_{l+1}-n_\infty\ge(B/A)(\mathfrak b_l-n_\infty)\), so
\(\mathfrak b_l\ge n_\infty+(B/A)^l(n-n_\infty)\ge n_\infty(1+(B/A)^l)\).
As \(B>A\), we have \(n_\infty>A\), so every step can be applied. Moreover
\(\mathfrak b_{l+1}\ge(B/A)\mathfrak b_l(1-u_l)\) with
\(u_l=A/\mathfrak b_l\le(1-A/B)/(1+(B/A)^l)\), so \(u_l\le\frac12\)
and \(u_l\le(1-A/B)(A/B)^l\). Since \(-\log(1-u)\le2u\) for
\(0\le u\le\frac12\) and \(\sum_{l\ge0}(1-A/B)(A/B)^l=1\), we
get \(\log\mathfrak b_j\ge\log n+j\log(B/A)-2\).

\subsection{The Hankel lemmas}
\label{subsec:proof-hankel}

\begin{proof}[Proof of Lemma~\ref{lem:hankel-ranks}]
For \(C(r,b)\), the entry at row \((\zeta,a)\) and column \((\eta,e)\)
is \([\zeta=a+e+\eta]=\sum_{m=0}^{b+v-2}[\zeta=a+m][m=e+\eta]\).
So \(\mathcal L_u(C(r,b))\) is a product of two matrices with inner
dimension \(b+v-1\). Conversely, the rows \((m,0)\) and the columns
\((\min(m,b-1),m-\min(m,b-1))\), for \(0\le m\le b+v-2\), form an
identity submatrix; the rows exist because \(m\le b+r-2\), as
\(v\le r\).

For \(C^{\mathrm{mid}}(r,c)\), put \(a'=u-1-a\) and \(e'=v-1-e\). As
\(r-1=(u-1)+(v-1)\), the entry is
\([\eta=r-1-a-e+\zeta]=[\eta=\zeta+a'+e']
=\sum_{m=0}^{c+u-2}[m=\zeta+a'][\eta=m+e']\), a product
with inner dimension \(c+u-1\). The rows with
\((\zeta,a')=(\min(m,c-1),m-\min(m,c-1))\) and the columns
with \((\eta,e')=(m,0)\), for \(0\le m\le c+u-2\), form an
identity submatrix; the columns exist because \(u\le r\). All matrices
here have entries \(0\) and \(1\), so the ranks are the same over every
field.
\end{proof}

\begin{proof}[Proof of Lemma~\ref{lem:hankel-gain}]
If \(r\) is odd, take \(u=v=(r+1)/2\). Then
\(c+u-1=c+v-1=(r+2c-1)/2\), and~\eqref{eq:hankel-bound} gives
\(M/q\le(\hat b/\hat c)^N\) with \(\hat b=r+2b-1\ge\hat c=r+2c-1\ge2\).
Since \(\hat b^2(\hat c^2-1)\le\hat c^2(\hat b^2-1)\), we get
\(\hat b^2/\hat c^2\le\varphi_r(b)/\varphi_r(c)\). If \(r\) is even, write
\eqref{eq:hankel-bound} for \((u,v)=(\frac r2,\frac r2+1)\) and for
\((u,v)=(\frac r2+1,\frac r2)\), and put \(c_+=c+\frac r2\) and
\(c_-=c+\frac r2-1\ge1\). The right sides are
\(\sum_gc_+^{N-\nmid(g)}c_-^{\nmid(g)}\) and
\(\sum_gc_-^{N-\nmid(g)}c_+^{\nmid(g)}\), and by the
Cauchy--Schwarz inequality their product is at least
\(M^2(c_+c_-)^N\). The product of the left sides is
\(q^2\bigl((b+\frac r2)(b+\frac r2-1)\bigr)^N\). Finally
\(4(y+\frac r2)(y+\frac r2-1)=\varphi_r(y)\).
\end{proof}

\begin{proof}[Proof of Lemma~\ref{lem:hankel-potential}]
(a) \(\partial_r\Psi=2r+4b-2>0\) and \(\partial_b\Psi=4r+2b-2>0\).

(b) Fix \(r\) and put \(\psi(y)=\frac23\log\Psi(r,y)-\frac12\log\varphi_r(y)\)
for real \(y\ge1\). With \(\hat y=r+2y-1\ge2\) and \(\delta=r-y\), we have
\(\varphi_r(y)=\hat y^2-1\) and \(\varphi_r'(y)=4\hat y\), and expanding
both sides shows
\[
 \tfrac23(4r+2y-2)\,\varphi_r(y)-2\hat y\,\Psi(r,y)=\tfrac13\bigl(2\hat y\delta^2-4\delta+\hat y\bigr).
\]
The left side is \(\Psi(r,y)\varphi_r(y)\psi'(y)\), and
\(\hat y(2\hat y\delta^2-4\delta+\hat y)=2(\hat y\delta-1)^2+\hat y^2-2>0\). As \(\Psi\) and
\(\varphi_r\) are positive, \(\psi\) is increasing, so
\(\psi(b)>\psi(c)\).

(c) With \(s=r+b-1\), we have \(\Psi=s^2-\frac76+2rb\) and
\(rb\le(s+1)^2/4\), so \(\Psi\le\frac32s^2+s-\frac23\), and
\(\frac{11}6s^2-(\frac32s^2+s-\frac23)=\frac13(s-1)(s-2)\ge0\) for
\(s\ge2\). By (a), \(\Psi(r,b)\ge\Psi(1,b)\), and
\((6\Psi(1,b))^2-121b^3=36(b-1)^4+167(b-1)^3+345(b-1)^2+165(b-1)\ge0\)
as \(b\ge1\).
\end{proof}

\subsection{Contraction and endpoints}
\label{subsec:lower-endpoint}

\begin{lemma}[Contraction]\label{lem:lower-contraction}
Every legal stage \((r,b)\to(c,r)\) of arity \(d\) satisfies
\(\hat\Psi(c,r)\le\mu_d^{-2}\,\hat\Psi(r,b)\).
\end{lemma}

\begin{proof}
Write \(\lambda=\lambda_d\) and \(\nu=\lambda+\frac1{d+2}\), so that
\(\lambda\nu=\lambda^2+\frac\lambda{d+2}=\frac d{d+2}\). Put
\(\tilde r=r-\frac13\), \(\tilde b=b-\frac13\) and
\(c^*=(db-r+1)/(d+2)\). Legality gives \(c\le c^*\), and
\(\tilde c^*=c^*-\frac13=(d\tilde b-\tilde r)/(d+2)\). Comparing the
coefficients of \(\tilde r^2\), \(\tilde r\tilde b\) and \(\tilde b^2\),
and using \(\lambda\nu=\frac d{d+2}\), shows
\[
 \lambda^2\hat\Psi(r,b)-\hat\Psi(c^*,r)
 =\lambda^2(\tilde r^2+4\tilde r\tilde b+\tilde b^2)
 -(\tilde c^{*2}+4\tilde c^*\tilde r+\tilde r^2)
 =(1-\nu^2)(\tilde r-\lambda\tilde b)^2 .
\]
For instance, the coefficients of \(\tilde b^2\) are
\(\lambda^2-\frac{d^2}{(d+2)^2}=\lambda^2(1-\nu^2)\). Since \(f_d\)
increases on \([0,\infty)\), \(f_d(\lambda)=0\) and
\(f_d(\frac{d+1}{d+2})=1\), we have \(\lambda<\frac{d+1}{d+2}\) and
\(0<\nu<1\). Hence \(\hat\Psi(c,r)\le\hat\Psi(c^*,r)\le\lambda^2\hat\Psi(r,b)\).
\end{proof}

\begin{lemma}[Endpoints]\label{lem:lower-endpoint}
All chain data satisfy
\(\log s\ge\log\kappa_0+\sum_{i=1}^J\log\mu_{d_i}+\frac45\log t\).
\end{lemma}

\begin{proof}
Let \((r,b)\to(c_i,r)\) be the stage of
Definition~\ref{def:class}(b). By Lemma~\ref{lem:lower-contraction},
symmetry and monotonicity,
\(\hat\Psi(r_{i+1},b_{i+1})\le\hat\Psi(c_i,r)\le\mu_{d_i}^{-2}\hat\Psi(r_i,b_i)\).
Also \(\hat\Psi(1,t)\le\hat\Psi(r_{J+1},b_{J+1})\). Hence
\(\hat\Psi(1,t)\le\hat\Psi(r_1,b_1)\prod_i\mu_{d_i}^{-2}\).

\emph{Start.} Since \(2xy\le\frac12(x+y)^2\), we have
\(x^2+4xy+y^2=(x+y)^2+2xy\le\frac32(x+y)^2\) for real \(x,y\). With \(x+y=r_1+b_1-\frac23=s+\frac13\), this gives
\(\hat\Psi(r_1,b_1)\le\frac32(s+\frac13)^2\le\frac94s^2\), because
\(12\bigl(\frac94s^2-\frac32(s+\frac13)^2\bigr)=9s^2-12s-2=(s-2)(9s+6)+10>0\)
for \(s\ge2\) (see the proof of Lemma~\ref{lem:recycling}(b)).

\emph{End.} \(\hat\Psi(1,t)=t^2+2t-\frac13\). We show
\(\hat\Psi(1,t)\ge\frac94\kappa_0^2t^{8/5}\), where
\(\frac94\kappa_0^2=\frac{3969}{1600}\). For \(t=1\) this is
\(\frac83\ge\frac{3969}{1600}\). For \(t\ge2\), we have
\(2t-\frac13\ge\frac{11}6t\), and the weighted arithmetic--geometric
mean inequality gives
\[
 t^2+\tfrac{11}6t=\tfrac35\cdot\tfrac53t^2+\tfrac25\cdot\tfrac{55}{12}t
 \ge\bigl(\tfrac53\bigr)^{3/5}\bigl(\tfrac{55}{12}\bigr)^{2/5}t^{8/5}.
\]
Here
\(\bigl(\frac53\bigr)^3\bigl(\frac{55}{12}\bigr)^2=\frac{378125}{3888}>97>94>\bigl(\frac{3969}{1600}\bigr)^5\).

Combining, \(\frac94\kappa_0^2t^{8/5}\le\frac94s^2\prod_i\mu_{d_i}^{-2}\).
Take half the logarithm.
\end{proof}

\begin{lemma}[Source and terminal]\label{lem:source-terminal}
All chain data satisfy:
\begin{enumerate}
\item[(a)] \(t\le s\);
\item[(b)] if \(d_i=d\) for some \(i\), then \(s\ge(1+\frac1d)t+\frac1d\);
\item[(c)] \(s\ge\frac{40}{49}Zt+\frac{76}{147}Z-\frac13\ge\frac{40}{49}Zt\),
where \(Z=\prod_i\mu_{d_i}\).
\end{enumerate}
\end{lemma}

\begin{proof}
(a) A legal stage \((r,b)\to(c,r)\) has \(c<b\), so the pair \((c,r)\)
has a smaller sum of entries than \((r,b)\) and no larger entry; crops
and exchanges increase neither. Hence
\(t\le\max(r_{J+1},b_{J+1})\le\max(r_1,b_1)\le s\).

(b) Let \((r,b)\to(c,r)\) be the stage of arity \(d\). By the proof
of~(a), \(s\ge r+b-1\) and \(t\le\max(c,r)\). If \(t\le r\), then
\(db\ge(d+2)c+r-1\ge r+d+1\), so \(s\ge(1+\frac1d)r+\frac1d\). If
\(t\le c\), then \(s\ge b\ge(1+\frac2d)c\ge(1+\frac1d)c+\frac1d\).

(c) The proof of Lemma~\ref{lem:lower-endpoint} gives
\(\hat\Psi(r_1,b_1)\ge Z^2\hat\Psi(1,t)\) and
\(\hat\Psi(r_1,b_1)\le\frac32(s+\frac13)^2\), so
\(s+\frac13\ge Z\sqrt{\frac23\hat\Psi(1,t)}\). Also
\(\frac23(t^2+2t-\frac13)-\bigl(\frac{40}{49}t+\frac{76}{147}\bigr)^2
=\frac{2(t-1)(t+1763)}{7203}\ge0\). This gives the first inequality, and
the second holds because \(Z\ge1\) and \(\frac{76}{147}>\frac13\).
\end{proof}

\subsection{The gap and the type counts}
\label{subsec:lower-gap}

\begin{lemma}[Gap]\label{lem:lower-gap}
For \(d\ge2\), \(g_d\le2/d\) and
\(\omega_d-\frac94\ge\frac1{24(d+2)^2}\ge\frac1{96d^2}\). In particular,
\(\log\mu_d>\frac34g_d\).
\end{lemma}

\begin{proof}
First, \(g_d=\log(1+\frac2d)\le\frac2d\). Put
\(z=(d/(d+2))^{1/4}\in(0,1)\), so that \(d+2=2/(1-z^4)\) and
\(g_d=-4\log z\). Substituting \(d=2z^4/(1-z^4)\) gives
\[
 (1-z^4)f_d(z^3)=2z^6+z^3(1-z^4)-2z^4=z^3(1-z)^2(1-z^2),
\]
so \(f_d(z^3)=z^3(1-z)^2/(1+z^2)>0\). Since \(f_d\) increases on
\([0,\infty)\) and \(f_d(\lambda_d)=0\), we get \(\lambda_d<z^3\), and
\[
 \omega_d-\tfrac94=\frac{-3\log\lambda_d+9\log z}{g_d}=\frac3{g_d}\log\frac{z^3}{\lambda_d}.
\]
By the mean value theorem, \(f_d(z^3)=f_d(z^3)-f_d(\lambda_d)\le(2d+5)(z^3-\lambda_d)\),
as \(f_d'(y)=2(d+2)y+1\le2d+5\) on \([0,1]\). With
\(\log y\ge1-1/y\),
\[
 \log\frac{z^3}{\lambda_d}\ge\frac{z^3-\lambda_d}{z^3}
 \ge\frac{(1-z)^2}{(1+z^2)(2d+5)} .
\]
Now \(1-z^4=(1-z)(1+z+z^2+z^3)\le4(1-z)\) gives
\(1-z\ge\frac1{2(d+2)}\). With \(1+z^2\le2\), \(2d+5\le\frac92d\) and
\(g_d\le\frac2d\), we get
\(\omega_d-\frac94\ge3\cdot\frac1{4(d+2)^2}\cdot\frac1{2\cdot\frac92d\cdot\frac2d}=\frac1{24(d+2)^2}\).
Finally \(d+2\le2d\), and \(\omega_d>\frac94\) is
\(\log\mu_d>\frac34g_d\).
\end{proof}

\begin{lemma}[Type counts]\label{lem:lower-types}
Let \(d\ge2\) and \(N\ge1\), and let \(M\) be the sector count of a type
of width \(N\) at arity \(d\). Then:
\begin{enumerate}
\item[(a)] \(1\le M<(d+2)^N\);
\item[(b)] if \(M>d^N\), then \(N\ge3\) and \(d\le2N\);
\item[(c)] if \(N\ge3\) and \(d\le2N\), then \(M\le(d+2)^N/2\);
\item[(d)] if \(N\ge2(d+2)\), then \(M<12(d+2)^{N+1}/N\).
\end{enumerate}
\end{lemma}

\begin{proof}
(a) is Lemma~\ref{lem:types}.

(b) For \(N=1\) the counts are \(d-1\), \(2\) and \(1\), all at most
\(d\). For \(N=2\) they are \((d-1)^2\), \(4(d-1)\), \(2(d-1)\), \(4\),
\(4\) and \(1\), all at most \(d^2\), since \(d^2-4(d-1)=(d-2)^2\). If
\(d\ge2N+1\), then for the type \((\nK,\nO,\nmid)\), using
\(N!/\nK!\le N^{\nO+\nmid}\),
\[
 \frac M{(d-1)^N}\le\frac1{\nO!\,\nmid!}\Bigl(\frac{2N}{d-1}\Bigr)^{\nO}
 \Bigl(\frac N{d-1}\Bigr)^{\nmid}\le1 ,
\]
so \(M\le(d-1)^N<d^N\).

(c) Put \(p_{\mathrm K}=\frac{d-1}{d+2}\), \(p_{\mathrm O}=\frac2{d+2}\)
and \(p_{\mathrm M}=\frac1{d+2}\). Then \(M/(d+2)^N\) is the probability
that \(N\) independent trials with outcome probabilities
\(p_{\mathrm K},p_{\mathrm O},p_{\mathrm M}\) produce the type. This is at
most \(\pi(\nK)\), where
\(\pi(a)=\binom Nap_{\mathrm K}^a(1-p_{\mathrm K})^{N-a}\) is the
probability of \(a\) kernel outcomes. We have \(\pi(a+1)\ge\pi(a)\) if and
only if \(a+1\le(N+1)p_{\mathrm K}\), and
\(1\le(N+1)p_{\mathrm K}\le N\), because \(N(d-1)\ge3\) and
\(d-1\le3N\). So \(\pi\) has a maximum at some \(a\) with
\(1\le a\le N-1\) (if \((N+1)p_{\mathrm K}=N\), then \(\pi(N-1)=\pi(N)\)).
At such an \(a\), the ratios \(\pi(a-1)/\pi(a)\) and \(\pi(a+1)/\pi(a)\)
have product \(\frac a{a+1}\cdot\frac{N-a}{N-a+1}\ge\frac14\), hence sum
at least \(1\). So \(2\pi(a)\le\pi(a-1)+\pi(a)+\pi(a+1)\le1\).

(d) Let \(\tau=(\tau_{\mathrm K},\tau_{\mathrm O},\tau_{\mathrm M})\) be a
type that maximizes the probability of (c), and let \(P_{\max}\) be this
probability; then \(M/(d+2)^N\le P_{\max}\). Moving one trial from class
\(j\) to class \(i\) multiplies the probability by
\(\tau_jp_i/((\tau_i+1)p_j)\le1\), and \(\tau_jp_i\le(\tau_i+1)p_j\) holds
trivially if \(\tau_j=0\). Summing over \(j\ne i\) gives
\(p_i(N-\tau_i)\le(\tau_i+1)(1-p_i)\), that is,
\(\tau_i\ge(N+1)p_i-1\). Since \(N\ge2d+4\), this gives
\(\tau_{\mathrm K}\ge\frac{N+1}4-1\ge\frac N8\),
\(\tau_{\mathrm O}\ge\frac N{d+2}\) and
\(\tau_{\mathrm M}\ge\frac N{2(d+2)}\ge1\). For integers \(y\ge1\),
\[
 \sqrt y\,(y/e)^y\le y!\le e\sqrt y\,(y/e)^y .
\]
Indeed, \(\log\) is concave, so the midpoint rule gives
\(\log y!\ge\int_{1/2}^{y+1/2}\log u\,du\ge y\log y-y+\frac12\log y\), and
the trapezoid rule gives
\(\log y!-\frac12\log y\le\int_1^y\log u\,du=y\log y-y+1\). Since
\(\sum_i\tau_i=N\), these bounds and \(\log u\le u-1\) give
\[
 P_{\max}\le e\sqrt{\frac N{\tau_{\mathrm K}\tau_{\mathrm O}\tau_{\mathrm M}}}\,
 \prod_i\Bigl(\frac{Np_i}{\tau_i}\Bigr)^{\tau_i}
 \le e\sqrt{\frac N{\tau_{\mathrm K}\tau_{\mathrm O}\tau_{\mathrm M}}}
 \le\frac{4e(d+2)}N<\frac{12(d+2)}N .\qedhere
\]
\end{proof}

\subsection{The resource lower bound}
\label{subsec:lower-main-proof}

\begin{proof}[Proof of Theorem~\ref{thm:lower}]
Put \(\Lambda_1=\log(1/\eps)\), \(d_\eps=(96\eps)^{-1/2}\) and
\(p_0=\frac94+\eps\). The hypothesis \(\eps<1/9216\) enters through four
conditions that it implies: \(\eps\le\frac3{20}\), \(\eps\le\frac9{64}\),
\(\eps\le\frac9{16384}\) and \(\eps<\frac1{9216}\) itself.

\emph{Step 0: reduction.} Let \(p'\le p_0\) be the pooled exponent of
the given chain scheme. By Lemma~\ref{lem:recycling}(b), the same chain
data with the source supply give a chain scheme
\(r_*\MM_{\bar A}\restr m_*\MM_{\bar B}\) with pooled exponent
\(p\le p'\le p_0\). As \(\eps\le\frac3{20}\),
\(p_0\le\frac{12}5<3\). The quantities \(r_*\), \(m_*\), \(W\), \(s\),
\(t\) and \(d_i\) do not depend on the supply.

\emph{Step 1: normalization.} Put \(w_{ij}=N_{ij}k_{ij}/W\), so that
\(\sum_jw_{ij}=1\), and \(\delta_{ij}=N_{ij}\log(d_i+2)-\log M_{ij}\),
which is positive by Lemma~\ref{lem:lower-types}(a). Taking logarithms
in~\eqref{eq:class-counts} and dividing by \(2W\) gives
\[
 \frac{\log(\bar B/\bar A)}{2W}=\log t+\sum_{i,j}w_{ij}\Bigl(g_{d_i}-\frac{\delta_{ij}}{N_{ij}}\Bigr),
 \qquad
 \frac{\log(r_*/m_*)}{2W}\ge3\log s+\sum_{i,j}w_{ij}\frac{\log E_{ij}}{2N_{ij}k_{ij}},
\]
where the factors \((2M_{ij}-1)/M_{ij}\ge1\) were dropped.

\emph{Step 2: the profit inequality.} Since \(p\le p_0\),
\(\log(r_*/m_*)\le p_0\log(\bar B/\bar A)\). By
Lemma~\ref{lem:lower-endpoint}, \(3\log\mu_d=\omega_dg_d\) and
\(\sum_jw_{ij}=1\), we have
\(3\log s\ge3\log\kappa_0+\sum_{i,j}w_{ij}\omega_{d_i}g_{d_i}+\frac{12}5\log t\).
Inserting both into Step~1 gives
\[
 \sum_{i,j}w_{ij}\vartheta_{ij}\ge3\log\kappa_0+\bigl(\tfrac{12}5-p_0\bigr)\log t\ge3\log\kappa_0,
 \qquad
 \vartheta_{ij}=(p_0-\omega_{d_i})g_{d_i}-\frac{p_0\delta_{ij}}{N_{ij}}-\frac{\log E_{ij}}{2N_{ij}k_{ij}} .
\]
Since \(\omega_{d_i}>\frac94\) (Lemma~\ref{lem:lower-gap}),
\(\delta_{ij}>0\) and \(E_{ij}\ge1\), every block has
\begin{equation}\label{eq:lower-theta}
 \vartheta_{ij}<(p_0-\omega_{d_i})g_{d_i}<\eps g_{d_i}.
\end{equation}
Let \(I\) be the set of blocks with \(\vartheta_{ij}>0\). Then
\(\sum_{(i,j)\in I}w_{ij}\vartheta_{ij}\ge3\log\kappa_0>0\), so \(I\) is not
empty.

\emph{Step 3: a profitable block.} Fix \((i,j)\in I\), and write \(d\),
\(N\), \(k\), \(M\), \(\delta\), \(E\), \(g\) and \(\omega\) for its
quantities.
\begin{itemize}
\item By~\eqref{eq:lower-theta}, \(\omega<p_0\). By
Lemma~\ref{lem:lower-gap}, \(\frac1{96d^2}\le\omega-\frac94<\eps\), so
\(d>d_\eps\).
\item As \(\vartheta>0\) and \(p_0-\omega<\eps\),
\begin{equation}\label{eq:lower-star}
 \frac{p_0\delta}N<\eps g\qquad\text{and}\qquad\frac{\log E}{2Nk}<\eps g .
\end{equation}
\item Hence \(\frac1N\log M=\log(d+2)-\frac\delta N>\log(d+2)-g=\log d\),
since \(\eps<p_0\). So \(M>d^N\).
\item By Lemma~\ref{lem:lower-types}(b) and~(c), \(\delta\ge\log2\). By
\eqref{eq:lower-star} and \(g\le\frac2d\),
\(N>\frac{p_0\delta}{\eps g}\ge\frac94\log2\cdot\frac d{2\eps}>\frac{9d}{16\eps}\).
\item As \(\eps\le\frac9{64}\), \(N>4d\ge2(d+2)\), so
Lemma~\ref{lem:lower-types}(d) applies. With \(d+2\le2d\), it gives
\(\delta>\log\frac N{24d}>\log\frac3{128\eps}\ge\frac{\Lambda_1}2\), the last step
because \(\eps\le\frac9{16384}\).
\item By~\eqref{eq:lower-star} again,
\(N>\frac{p_0\delta}{\eps g}>\frac94\cdot\frac{\Lambda_1}2\cdot\frac d{2\eps}=\frac{9d\Lambda_1}{16\eps}\).
\item \(E=M+(6k-1)(M-1)\ge M>d^N\), so~\eqref{eq:lower-star} gives
\(\frac{\log d}{2k}<\eps g\le\frac{2\eps}d\), that is,
\(k>\frac{d\log d}{4\eps}\).
\item As \(\eps<\frac1{9216}\), \(\log d_\eps=\frac12(\Lambda_1-\log96)>\frac{\Lambda_1}4\).
Since \(d>d_\eps>1\) and \(y\log y\) increases for \(y\ge1\),
\(d\log d>d_\eps\log d_\eps>\frac{d_\eps\Lambda_1}4\).
\end{itemize}
These are the bounds on \(d_i\), \(N_{ij}\) and \(k_{ij}\) in the theorem.
Since \(W\ge Nk\) and \(d_\eps^2=\frac1{96\eps}\),
\[
 W>\frac{9d_\eps\Lambda_1}{16\eps}\cdot\frac{d_\eps\Lambda_1}{16\eps}=\frac{3\Lambda_1^2}{8192\,\eps^3}.
\]

\emph{Step 4: the profitable sum.} For \((i,j)\in I\),
\eqref{eq:lower-theta} and Lemma~\ref{lem:lower-gap} give
\(\vartheta_{ij}<\eps g_{d_i}\le\frac{2\eps}{d_i}<\frac{2\eps\log d_i}{d_\eps\log d_\eps}\),
as \(d_i\log d_i>d_\eps\log d_\eps\). Weight by \(w_{ij}\) and sum over
\(I\):
\[
 \Sigma_I=\sum_{(i,j)\in I}w_{ij}\log d_i
 >\frac{3\log\kappa_0\,d_\eps\log d_\eps}{2\eps}>\frac{3\log\kappa_0\,d_\eps\Lambda_1}{8\eps}.
\]

\emph{Step 5: conclusion.} Since \(\sum_jw_{ij}=1\), Steps~3 and~4 give
\[
 \log\bar A=2W\sum_{i,j}w_{ij}\log d_i\ge2W\,\Sigma_I
 >2\cdot\frac{3\Lambda_1^2}{8192\,\eps^3}\cdot\frac{3\log\kappa_0\,d_\eps\Lambda_1}{8\eps}
 =\frac{9\log\kappa_0}{32768\sqrt{96}}\,\eps^{-9/2}\Lambda_1^3 .
\]
As \(\log(1+y)\ge y/(1+y)\), \(\log\kappa_0\ge\frac1{21}\), and
\(\sqrt{96}<10\), so the coefficient exceeds
\(\frac9{21\cdot32768\cdot10}>1.3\cdot10^{-6}\). Under the source
supply, \(A=\bar A\), which proves~(c).

For~(a), the proof of Lemma~\ref{lem:recycling}(b) gives \(r_*/m_*\ge s^{6W}>1\).
Since \(p\le3\),
\(\log(\bar B/\bar A)\ge\frac13\log(r_*/m_*)\ge2W\log s\). By
\eqref{eq:class-counts}, \(\bar B=t^{2W}m_*^{1/3}\), so
\(\frac13\log m_*\ge\log\bar A+2W\log(s/t)\). By
Lemma~\ref{lem:source-terminal}(a), \(t\le s\), so
\(\log m_*\ge3\log\bar A\). Finally
\(3\cdot\frac9{21\cdot32768\cdot10}>3.9\cdot10^{-6}\).

For~(b), Step~2 and~\eqref{eq:lower-theta} give
\(3\log\kappa_0\le\sum_{i,j}w_{ij}\vartheta_{ij}<\eps\sum_ig_{d_i}\). By
Lemmas~\ref{lem:lower-endpoint} and~\ref{lem:lower-gap}, with
\(\log\kappa_0>0\) and \(t\ge1\),
\(\log s\ge\sum_i\log\mu_{d_i}>\frac34\sum_ig_{d_i}>\frac{9\log\kappa_0}{4\eps}\).
For the given scheme, \(p'\le p_0\le\frac{12}5\) gives
\[
 \log(B/A)=\frac{\log(r_*/m_*)}{p'}\ge\frac{\log(r_*/m_*)}{p_0}
 \ge\frac{6W\log s}{12/5}>\frac{45\log\kappa_0}8\cdot\frac W\eps
 >\frac{45}{8\cdot21}\cdot\frac{3\Lambda_1^2}{8192\,\eps^4},
\]
and \(\frac{135}{8\cdot21\cdot8192}=\frac{45}{458752}>9.8\cdot10^{-5}\).
\end{proof}

\begin{remark}[Joint exactification of one stage]\label{rem:joint-exactification}
Suppose that the blocks of each stage are exactified together, so that
the factors \(E_{ij}\) of \(r_*\) are replaced by
\(E_i=1+\sum_j6k_{ij}(M_{ij}-1)\). Then Theorem~\ref{thm:lower} still
holds, except for the bound on \(k_{ij}\). In the proof above, Step~1
holds with \(\sum_i\log E_i/(2W)\) in place of
\(\sum_{i,j}w_{ij}\log E_{ij}/(2N_{ij}k_{ij})\), and Step~2 gives
\(\sum_i\vartheta_i\ge3\log\kappa_0\) with one profit term per stage,
\[
 \vartheta_i=\sum_jw_{ij}\Bigl((p_0-\omega_{d_i})g_{d_i}-\frac{p_0\delta_{ij}}{N_{ij}}\Bigr)-\frac{\log E_i}{2W}
 <\eps g_{d_i}.
\]
If \(\vartheta_i>0\), then some block \(j\) of stage \(i\) has
\((p_0-\omega_{d_i})g_{d_i}>p_0\delta_{ij}/N_{ij}\), and Step~3 gives the
bounds on \(d_i\) and \(N_{ij}\) for this block. Moreover
\(\log E_i/(2W)<\eps g_{d_i}\le2\eps/d_i\) and
\(E_i\ge M_{ij}>d_i^{N_{ij}}\) give \(W>N_{ij}d_i\log d_i/(4\eps)\),
which replaces \(W\ge N_{ij}k_{ij}\) in the bound on \(W\). Steps~4
and~5 run over the stages with \(\vartheta_i>0\), with \(\Sigma_I\)
replaced by the sum of \(\log d_i\) over these stages.
\end{remark}

\subsection{The recycled main construction}
\label{subsec:lower-recycled-proof}

\begin{proof}[Proof of Corollary~\ref{cor:recycled-main}]
We use the notation of Theorem~\ref{thm:main} and its proof. Here
\(A\) and \(B\) are the sides of Theorem~\ref{thm:main}, that is, of the
source supply, \(W=Nk\), \(t=1\), every stage has one block with the
largest sector count \(M_i\), and, as there, \(a_N=2\log(N+1)/N\) and
\(\Lambda=\log(4096/\eps)\). Read forwards, the chain has \(L\ge1\) bulk
stages of arity \(D\), followed by ramp stages, and every arity
satisfies \(d_i\le24D-3\) (Lemma~\ref{lem:defect}). We use these facts
from the proof of Theorem~\ref{thm:main}: \(N=k\ge256DH/\eps\ge4608/\eps\)
(as \(D\ge2\) and \(H\ge9\)), \(a_N<3\eps/(128D)\) and
\(Lg_D\ge64/\eps\) (Steps~0 and~1);
\(\log(B/A)\ge\frac{497}{512}\cdot2NkG\) with \(G\ge64/\eps\)
(Steps~1 and~4; see also~\eqref{eq:seeds-gain}); and \(D\le3\eps^{-1/2}\),
\(J\le240\,\eps^{-3/2}\) and \(N\le2305\,\eps^{-3/2}\Lambda\) (Step~5).

\emph{Funding.} By Lemma~\ref{lem:types},
\(\log M_i\ge N\log(d_i+2)-2\log(N+1)\), so
\(\gamma_i=\log(V_i/U_i)=2k\log M_i-2Nk\log d_i\ge2Nk(g_{d_i}-a_N)\).
\begin{itemize}
\item Every stage gains: \(g_{d_i}\ge\frac2{d_i+2}>\frac1{12D}\), as
\(d_i+2<24D\), and \(a_N<\frac{3\eps}{128D}\le\frac3{128D}<\frac1{12D}\).
So \(\gamma_i>0\).
\item The bulk pays for the larger arities of the ramp:
\(La_N<(\frac{64D}\eps+1)\frac{3\eps}{128D}\le\frac32+\frac3{256}<2\), so
\(\sum_{i\le L}\gamma_i\ge2Nk(Lg_D-La_N)>2Nk(\frac{64}\eps-2)\ge124Nk\).
\end{itemize}
For \(i\le L\), \(U_i=D^{2Nk}\) and \(\prod_{l<i}U_l/V_l\le1\), with
equality for \(i=1\). For \(i>L\), \(U_i\le(24D)^{2Nk}\) and
\(\prod_{l<i}U_l/V_l<e^{-124Nk}<24^{-2Nk}\). Hence
\(\max_iU_i\prod_{l<i}U_l/V_l=D^{2Nk}\), and
Lemma~\ref{lem:recycling}(c) applies to \(A'=4JD^{2Nk}\): the recycling
supply is a supply, and \(B'/A'\ge B/(2A)\). By
Lemma~\ref{lem:recycling}(a), \(B'/A'\le B/A\).

\emph{Exponent.} We have
\(\log(B/A)\ge\frac{497}{512}\cdot\frac{128}\eps Nk>\frac{124Nk}\eps\ge\frac{124\cdot4608^2}\eps\),
so \(\log(B'/A')\ge\log(B/A)-\log2>0\), and \(B'>A'\). The two pooled
exponents have the same numerator \(p\log(B/A)\), so
\[
 p'-p\le\frac{p\log2}{\log(B/A)-\log2}
 <\frac{\frac{11}4\log2}{124\cdot4608^2-1}\,\eps<10^{-9}\eps,
\]
using \(p<\frac{11}4\) and \(\eps\le1\). With
\(p<\frac94+\frac{719}{1988}\eps\) (Theorem~\ref{thm:main}) and
\(\frac{719}{1988}+10^{-9}<\frac12\), this gives
\(p'<\frac94+\frac\eps2\). Also \(p'\ge p>\frac94\), by
Lemma~\ref{lem:recycling}(b) and Theorem~\ref{thm:main}.

\emph{Source side.} \(\log A'=\log(4J)+2Nk\log D\). Here
\(4J\le960\,\eps^{-3/2}\le(4096/\eps)^{3/2}\), so
\(\log(4J)\le\frac32\Lambda\); \(2\log D\le\log(9/\eps)\le\Lambda\); and
\(Nk=N^2\le2305^2\eps^{-3}\Lambda^2\). As \(\Lambda\ge\log4096>8.3\),
\(\frac32\Lambda<0.022\,\Lambda^3\). Hence
\[
 \log A'<(2305^2+0.022)\,\eps^{-3}\Lambda^3<5.32\cdot10^6\,\eps^{-3}\Lambda^3 .
\]

\emph{Target side.} Since \(M_i<(d_i+2)^N\) and \(t=1\),
\(\log(B/A)<2NkG\), where \(G=\log F(R,R)+Lg_D\)
(Lemma~\ref{lem:defect}). Here \(Lg_D\le(\frac{64D}\eps+1)\frac2D\le\frac{128}\eps+1\).
By~\citep[(3.6)]{FinCon26}, \(\log F(R,R)\le\frac43\log(2R-1)\), and
\(2R-1<16D\le48\,\eps^{-1/2}\); since \(48^4<4096^2\), this gives
\(\log F(R,R)<\frac23\Lambda\). The function \(\eps\log(4096/\eps)\)
increases on \((0,1]\), so \(\eps\Lambda\le\log4096<8.32\). Hence
\[
 \log(B/A)<Nk\bigl(\tfrac{256}\eps+2+\tfrac43\Lambda\bigr)
 \le2305^2\eps^{-4}\Lambda^2\bigl(256+2\eps+\tfrac43\eps\Lambda\bigr)
 <1.43\cdot10^9\,\eps^{-4}\Lambda^2 .
\]
By Lemma~\ref{lem:recycling}(a), \(\log B'\le\log A'+\log(B/A)\), and
\(\log A'<5.32\cdot10^6\,\eps^{-3}\Lambda^3<4.5\cdot10^7\,\eps^{-4}\Lambda^2\),
as \(\eps\Lambda<8.32\). So
\(\log B'<1.5\cdot10^9\,\eps^{-4}\Lambda^2\).
The restriction holds over \(\Q\) by Lemma~\ref{lem:recycling}, with the
\(r_*\) and \(m_*\) of Theorem~\ref{thm:main}.
\end{proof}

\subsection{Small batches}
\label{subsec:lower-small-proof}

\begin{proof}[Proof of Proposition~\ref{prop:small-sources}]
(a) For integers \(N,k\ge1\), \(d\ge2\) and \(M\ge1\), put
\(E=1+6k(M-1)\) and
\(\upsilon(N,d,k,M)=\frac1{2Nk}\log\bigl(d^{6Nk}(2M-1)^{6k}E/M^{12k}\bigr)\).
Let \(\upsilon_{ij}=\upsilon(N_{ij},d_i,k_{ij},M_{ij})\) and let
\(\bar\upsilon_i=\sum_j(N_{ij}k_{ij}/W)\upsilon_{ij}\), a weighted average. Under
the source supply, \(\log A=\sum_{i,j}2N_{ij}k_{ij}\log d_i\), so
\eqref{eq:class-counts} gives
\begin{equation}\label{eq:small-margin}
 \frac{\log(r_*/m_*)-3\log(B/A)}{2W}=3\log\frac st+\sum_i\bar\upsilon_i ,
\end{equation}
and \(p<3\) if and only if this is negative. If \(M\le d^N\), then
\(\upsilon(N,d,k,M)\ge0\), because \(d^{6Nk}\ge M^{6k}\), \(2M-1\ge M\) and
\(E\ge1\). If a type of width \(N\) has sector count \(M>d^N\), then
\(N\ge3\) and \(d\le2N\) (Lemma~\ref{lem:lower-types}(b)).

Suppose that \(p<3\) and \(m_*<M_0^{12}\), and let \(\mathcal K\) be the
set of stages with \(\bar\upsilon_i<0\). For \(i\in\mathcal K\), let
\(N,k,M\) be the width, power and sector count of a block of stage \(i\)
with the least \(\upsilon_{ij}\), and \(\upsilon_i^\circ\) its \(\upsilon_{ij}\). Then
\(\upsilon_i^\circ\le\bar\upsilon_i<0\), so \(M>d_i^N\), \(N\ge3\), \(d_i\le2N\),
and \(M^{6k}\le m_*<M_0^{12}\). By~\eqref{eq:small-margin},
\(0>3\log(s/t)+\sum_{i\in\mathcal K}\upsilon_i^\circ\). If
\(\mathcal K=\emptyset\), this contradicts
Lemma~\ref{lem:source-terminal}(a). If \(\mathcal K=\{i\}\) and
\(d=d_i\), Lemma~\ref{lem:source-terminal}(b) gives \(s>(1+\frac1d)t\),
so
\[
 0>3\log\frac{d+1}d+\upsilon_i^\circ
 =\frac1{2Nk}\log\frac{(d+1)^{6Nk}(2M-1)^{6k}E}{M^{12k}},
\]
which needs \(M>(d+1)^N\) (as above) and contradicts check~(i) below. If
\(|\mathcal K|\ge2\), Lemma~\ref{lem:source-terminal}(c) and
\(\mu_d>1\) give
\(3\log(s/t)\ge3\log\frac{40}{49}+3\sum_{i\in\mathcal K}\log\mu_{d_i}\),
and check~(ii) gives \(3\log\mu_{d_i}+\upsilon_i^\circ\ge\log\frac{34}{25}\).
So \(0>3\log\frac{40}{49}+2\log\frac{34}{25}\), which is false, as
\(34^2\cdot40^3=73\,984\,000>73\,530\,625=25^2\cdot49^3\).

The two checks concern all \((N,d,k,M)\) with \(N\ge3\),
\(2\le d\le2N\), \(k\ge1\), \(M\) a sector count of width \(N\) at arity
\(d\), and \(M^{6k}<M_0^{12}\): (i) if \(M>(d+1)^N\), then
\((d+1)^{6Nk}(2M-1)^{6k}E\ge M^{12k}\); (ii) if \(M>d^N\), then
\(25^{2Nk}(vd)^{6Nk}(2M-1)^{6k}E\ge34^{2Nk}u^{6Nk}M^{12k}\), where
\(u/v\ge\lambda_d\) is a rational number fixed for each \(d\). Since
\(\log\mu_d\ge\log(v/u)\), (ii) gives the bound used above. Both lists
are finite, since \(d^{6Nk}<M^{6k}<2^{495}\) gives \(N\le82\), and
\(d\ge33\) would give \(N\ge17\) and \(d^{6N}\ge33^{102}>2^{495}\). The
supplementary script \path|shifted_chains_check_small_sources.py|
(Part~C) runs both checks in exact integer arithmetic, with \(u/v\) the
upper end of an enclosure of \(\lambda_d\) from an integer square root,
certified by \(f_d(u/v)\ge0\). All \(5\,792\) cases of (i) and
\(31\,677\) cases of (ii) pass.

(b) Every restriction \(r\MM_A\restr m\MM_B\) in the construction has
\(r\ge m\) (for chain schemes since \(2M_{ij}-1\ge M_{ij}\)), and the
operations preserve this. By induction on the construction, \(B>A\) and
\(r/m<(B/A)^3\) imply \(m\ge M_0^{12}\). For chain schemes this is (a);
classical restrictions have \(r/m=(B/A)^3\); and (i) keeps \(r\) and
\(m\) and does not increase \(B/A\). Under (ii) and (iii),
\(m=m^{(1)}m^{(2)}\), \(r/m=(r^{(1)}/m^{(1)})(r^{(2)}/m^{(2)})\) and
\(B/A\le(B^{(1)}/A^{(1)})(B^{(2)}/A^{(2)})\), since
\(f\le B^{(1)}/A^{(2)}\). So \(r/m<(B/A)^3\) gives
\(r^{(l)}/m^{(l)}<(B^{(l)}/A^{(l)})^3\), hence \(B^{(l)}>A^{(l)}\) as
\(r^{(l)}\ge m^{(l)}\), for some \(l\), and \(m\ge m^{(l)}\ge M_0^{12}\).
The composition is a restriction by the construction of
Remark~\ref{rem:composition}, which does not use \(B^{(l)}>A^{(l)}\).
Finally, the maps of \(r\MM_A\restr m\MM_B\) replace each variable of
\(r\MM_A\) by a linear form in the \(mB^2\) variables of the same leg of
\(m\MM_B\), and each of these occurs in \(m\MM_B\), hence in some form.
As \(B\ge2\), this gives at least \(mB^2\ge4M_0^{12}>2^{496}\) nonzero
coefficients per leg.

(c) The stage is legal, since \(4\le2\cdot182-361+1\), and \(s=542\).
So \(A=2^{92}\), \(B=361^{92}M_0^4\),
\(r_*=542^{276}(2M_0-1)^{12}(12M_0-11)\) and \(m_*=M_0^{12}\). The script
checks \(r_*A^3<m_*B^3\), that is \(p<3\), in exact integer arithmetic,
and its enclosures give \(p\in[2.999998425472,\,2.999998425473]\).
\end{proof}
